% !TEX root = ../main.tex
%

\section{Investigating facilitation through SDG}
\label{sec:facilitation}

\subsection{LLM facilitators constantly intervene}
\label{sec:facilitation:interventions}

\begin{figure}[ht]
	\centering
	\includegraphics[width=0.45\linewidth]{intervention_count_model.png}
	\includegraphics[width=0.45\linewidth]{intervention_count_strategy.png}
	\Description{Lineplots showing the number of interventions by strategy and by model, for each discussion turn. All lines are close to the 100\% mark, except for the `No Instructions' strategy and the Qwen-32B model, which hover around the 80\% and 30\% marks respectively.}
	\caption{Percentage of cumulative facilitator interventions averaged across discussions, showing how often the facilitator chose to respond up to each point in the conversation. Results are grouped by model (right) and facilitation strategy (left; see~\S\ref{sec:facilitation:strategies}). LLM facilitators consistently and constantly intervene in discussions, with little overall variation between models and strategies. Error bars indicate 95\% confidence intervals.}
	\label{fig:facilitation-model}
\end{figure}

Fig.~\ref{fig:facilitation-model} shows that LLM facilitators tend to intervene at nearly every opportunity, despite being instructed to act only when necessary (see Appendix~\ref{ssec:appendix:moderation_strategies}). This behavior is consistent across time, model families, model sizes, and facilitation strategies. The only notable exception is the Qwen-32B model, which exhibits comparatively lower rates of intervention; however, this is likely due to the inability of the model to produce toxic speech (Fig.~\ref{fig:models}, right). This analysis is an instance of applying multiple evaluation criteria to separate genuine behavioral patterns from artifacts arising from limitations in the \ac{sdg} simulations (\S\ref{sec:methodology:evaluation}). Our results confirm that:

\begin{finding}
	\label{finding:intervention}
	LLM facilitators \emph{can not decide when to intervene, and do so constantly} even when explicitly instructed to only intervene when necessary.  
\end{finding}

The inability to keep silent has not been reported in relevant literature to our knowledge, and \emph{is an example of ``debugging'' problems with LLMs}--a core motivation of our work. Without the discovery of this limitation it is very plausible that researchers investigating LLM facilitation would have spent a significant amount of time and money to set up experiments with human participants and numerous LLM configurations as proposed in recent work~\cite{korre2025evaluation,small-polis-llm}; only to find that this limitation would render the rest of their findings inadmissible.

Furthermore, the fact that interventions remain constant throughout the discussion without toxicity being lowered may imply that facilitation is not effective, because effective facilitation needs targeted, non-repetitive interventions~\cite{chen-etal-2025-whow, park_et_al_2012_facilitation, schroeder-etal-2024-fora}. It is plausible that this finding is replicable to human discussions to at least some extent, since the inability to choose when to speak seems to be a general issue with LLMs (\S\ref{sec:methodology:management:turn-taking}) and is invariant towards instruction prompts, model family and size, conversational topics, and experimental conditions. This suggests a universal limitation of LLM facilitators, rather than an observation conditioned on synthetic data.


\subsection{Facilitative strategies have an observable effect on conversational dynamics}
\label{sec:facilitation:strategies}

We test two facilitation strategies used in real-life and compare them with two baselines. We summarize both documents using a proprietary LLM (Google Gemini~\cite{geminiteam2025}), since LLMs have been proven to work well in these tasks~\cite{ts2024}, manually inspected the output, and removed instructions not relevant for our scenario .\footnote{Note that the process of turning sometimes extensive documents into short prompts (Appendix~\ref{sec:appendix:prompts}), necessitated by open-source LLMs, is necessarily imperfect. We leave the optimal derivation of strategy prompts to future work.}

\begin{enumerate}[nosep,noitemsep]
	\item \textbf{No Facilitator}: A \emph{baseline} where no facilitator is present.
	
	\item \textbf{No Instructions}: A \emph{baseline} where a LLM facilitator is present, but is provided only with basic instructions (e.g., “You are a moderator, keep the discussion civil”). This approach is already being used in some platforms \cite{Tsai_Deliberate_Lab_Open-Source_2025}. 
	
	\item \textbf{Moderation Guidelines}: A \emph{real-life} strategy based on guidelines given to human facilitators of the `Regulation Room' platform \cite{Cornell_eRulemaking2017}. The instructions are suitable for online fora, where facilitators also engage in content moderation (removing comments, see \S\ref{sec:related:facilitation}, and their effectiveness must be balanced by their throughput (e.g., ``Stick to a maximum of two questions, use simple and clear language, deal with off-topic comments'').
	
	\item \textbf{Constructive Communications}: A \emph{real-life} strategy~\cite{dimitra-guide} which approaches facilitation from a more personalized and indirect angle, forbidding facilitators from directly providing opinions or directions (e.g., ``Do not make decisions, be a guide, provide explanations''). This makes the strategy ideal for deliberative environments.
\end{enumerate}

To assess the impact of the proposed strategies in simulated discussions, we employ linear regression models as an interpretable and computationally efficient approach for identifying statistically significant patterns in toxicity scores obtained from the Perspective API (\S\ref{sec:scalability:experimental}). Specifically, we use \acp{mem}, which account for the hierarchical structure of the data by grouping comments within discussions. This is important, as comments from the same discussion are likely to be correlated in both content and toxicity. Each model estimates coefficients that capture the average change in toxicity relative to a baseline (``average'') discussion, along with corresponding standard errors and p-values.  Negative toxicity coefficients indicate lower toxicity relative to the baseline, while positive coefficients indicate higher toxicity. The baseline is defined through reference levels of categorical variables (e.g., ``No trolls'' for the troll condition, or ``No Facilitator'' for facilitation strategies). To quantify how toxicity evolves over the course of a discussion, we include an ordinal variable, \textit{Discussion Turn}, indicating the position of each comment. This variable is mean-centered, such that a value of 0 corresponds to the midpoint of the discussion, following standard practice in linear regression. For example, the model presented in Table~\ref{tab:regression-overall} can be written as:
\[
\textit{Toxicity}_{ij} = \beta_0 + \beta_1 CC_{ij} + \beta_2 MG_{ij} + \beta_3 MI_{ij}
+ \beta_4 (CC_{ij} \times T_{ij}) + \beta_5 (MG_{ij} \times T_{ij}) + \beta_6 (MI_{ij} \times T_{ij})
+ u_j + \epsilon_{ij}
\]
\noindent where $\textit{Toxicity}_{ij}$ denotes the toxicity score of comment $i$ in discussion $j$, $CC$, $MG$, and $MI$ are binary indicators for the presence of each facilitation strategy (coded as 1), with “No Facilitator” as the reference category (coded as 0), $T_{ij}$ represents the centered discussion turn., $\beta_0$ corresponds to the expected toxicity under the reference condition at the midpoint of the discussion,
$\beta_1 \ldots \beta_3$ estimate the average effect of each facilitation strategy, while $\beta_4 \ldots \beta_6$ capture how these effects vary over the course of the discussion via interaction with discussion turn. The term $u_j$ represents a discussion-level random effect, quantifying unobserved heterogeneity across discussions, and $\epsilon_{ij}$ is the residual error at the comment level.


The choice of facilitation strategy significantly shapes the dynamics of simulated discussions, both in terms of overall toxicity levels and how these evolve over time. As shown in Table~\ref{tab:regression-overall}, discussions without facilitation start from a baseline toxicity of 0.184. All three facilitation strategies--Constructive Communications (CC), Moderation Guidelines (MG), and Minimal Instructions (MI)--lead to significantly lower toxicity levels, with reductions of -0.040, -0.032, and -0.030, respectively. This indicates that any form of facilitation has a dampening effect on participant toxicity compared to having no facilitator at all. However, toxicity also evolves over time: the positive coefficient for discussion turn (T = 0.001) suggests a slight increase in toxicity as discussions progress under the baseline condition. Importantly, this trend is moderated differently depending on the facilitation strategy; for example, CC not only reduces baseline toxicity but also counteracts the tendency for toxicity to increase over time ($CC\times T = -0.002$), and MI shows a weaker but still significant mitigating effect ($MI\times T = -0.001$), while MG’s interaction is not statistically significant, suggesting it reduces initial toxicity but does not meaningfully alter its trajectory.


\begin{table}[ht]
	\centering
	\caption{\ac{mem} on LLM participant toxicity depending on the facilitator's strategy. Reference level is \texttt{No Facilitator}. Acronyms: CC = Constructive Communications, MG = Moderation Guidelines, MI = Minimal Instructions, T = Discussion Turn (centered number of comments in the discussion so far). $^*p<0.05$, $^{**}p<0.01$, $^{***}p<0.001$. Standard errors in parentheses.}
	\label{tab:regression-overall}
	\begin{tabular}{lcccccccc}
		\toprule
		& Intercept & CC & MG & MI & T & CC $\times$ T & MG $\times$ T & MI $\times$ T \\
		\midrule
		& \begin{tabular}[c]{@{}c@{}}$0.184^{***}$ \\ $(0.008)$\end{tabular}
		& \begin{tabular}[c]{@{}c@{}}$-0.040^{***}$ \\ $(0.011)$\end{tabular}
		& \begin{tabular}[c]{@{}c@{}}$-0.032^{**}$ \\ $(0.011)$\end{tabular}
		& \begin{tabular}[c]{@{}c@{}}$-0.030^{**}$ \\ $(0.011)$\end{tabular}
		& \begin{tabular}[c]{@{}c@{}}$0.001^{*}$ \\ $(0.000)$\end{tabular}
		& \begin{tabular}[c]{@{}c@{}}$-0.002^{**}$ \\ $(0.001)$\end{tabular}
		& \begin{tabular}[c]{@{}c@{}}$-0.001$ \\ $(0.001)$\end{tabular}
		& \begin{tabular}[c]{@{}c@{}}$-0.001^{*}$ \\ $(0.001)$\end{tabular} \\
		\bottomrule
	\end{tabular}
\end{table}


Table~\ref{tab:regression-facilitator} shows that the facilitation strategy also significantly affects the behavior of the LLM facilitator itself. In the absence of instructions, facilitator toxicity is already quite low (0.061). Both CC and MG further reduce this level by -0.020 and -0.023, respectively, suggesting that structured strategies encourage more neutral or restrained interventions. The effect of message order is effectively zero, and none of the interaction terms with time are significant, indicating that facilitator behavior remains stable throughout the discussion regardless of strategy. Given the already low baseline, these reductions are small in absolute terms but consistent, reinforcing the interpretation that different strategies primarily modulate how assertive the facilitator is.

\begin{table}[ht]
	\centering
	\caption{\ac{mem} on toxicity exhibited by the \emph{LLM facilitator} depending on the facilitation strategy. Reference level is \texttt{No Instructions}. Acronyms: CC = Constructive Communications, MG = Moderation Guidelines, T = message order (centered). $^*p<0.05$, $^{**}p<0.01$, $^{***}p<0.001$. Standard errors in parentheses.}
	\label{tab:regression-facilitator}
	\begin{tabular}{lcccccc}
		\toprule
		& Intercept & CC & MG & T & CC $\times$ T & MG $\times$ T \\
		\midrule
		Facilitator 	Toxicity 
		& \begin{tabular}[c]{@{}c@{}}$0.061^{***}$ \\ $(0.002)$\end{tabular}
		& \begin{tabular}[c]{@{}c@{}}$-0.020^{***}$ \\ $(0.003)$\end{tabular}
		& \begin{tabular}[c]{@{}c@{}}$-0.023^{***}$ \\ $(0.003)$\end{tabular}
		& \begin{tabular}[c]{@{}c@{}}$0.000$ \\ $(0.000)$\end{tabular}
		& \begin{tabular}[c]{@{}c@{}}$-0.000$ \\ $(0.000)$\end{tabular}
		& \begin{tabular}[c]{@{}c@{}}$-0.000$ \\ $(0.000)$\end{tabular} \\
		\bottomrule
	\end{tabular}
\end{table}

Finally, Table~\ref{tab:regression-response} examines whether facilitation strategies influence how individual users respond after being intervened upon. The strongest predictors of post-intervention toxicity are whether the participant is a troll (0.150) and their pre-intervention toxicity (0.172). In contrast, neither CC (-0.007) nor ER (-0.002) has a statistically significant direct effect on post-intervention toxicity, and their interactions with trolls are also non-significant. This suggests that facilitation strategies do not meaningfully alter the immediate behavior of the targeted user. Instead, the earlier findings imply that their impact operates at the discussion level--changing the broader conversational environment--rather than through direct behavioral correction of individual participants following an intervention. Taken together, these findings suggest that:

\begin{finding}
	\label{finding:strategy}
	Supplying LLM facilitators with different facilitation strategies \emph{shifts the overall trajectory of discussions}, as opposed to only having an immediate effect on the user who triggered an intervention. 
\end{finding}

\begin{table}[ht]
	\centering
	\caption{\ac{mem} predicting post-intervention toxicity. In a chain of user A commenting, followed by an intervention, followed by user A commenting again, we consider the first and last comments to be the pre- and post-intervention toxicity levels respectively. We control for pre-toxicity, since it should be highly correlated with post-toxicity. Reference levels: \texttt{No Instructions} for strategy, \texttt{False} for trolls. Acronyms: CC = Constructive Communications, ER = E-Rulemaking, TR = whether the participant is a troll, PI= Pre-Intervention toxicity. $^*p<0.05$, $^{**}p<0.01$, $^{***}p<0.001$. Standard errors in parentheses.}
	\label{tab:regression-response}
	\begin{tabular}{lccccccc}
		\toprule
		& Intercept & CC & ER & TR & CC $\times$ TR & ER $\times$ TR & PI \\
		\midrule
		Post-Interv. Toxicity
		& \begin{tabular}[c]{@{}c@{}}$0.084^{***}$ \\ $(0.006)$\end{tabular}
		& \begin{tabular}[c]{@{}c@{}}$-0.007$ \\ $(0.009)$\end{tabular}
		& \begin{tabular}[c]{@{}c@{}}$-0.002$ \\ $(0.009)$\end{tabular}
		& \begin{tabular}[c]{@{}c@{}}$0.150^{***}$ \\ $(0.007)$\end{tabular}
		& \begin{tabular}[c]{@{}c@{}}$-0.006$ \\ $(0.008)$\end{tabular}
		& \begin{tabular}[c]{@{}c@{}}$0.007$ \\ $(0.008)$\end{tabular}
		& \begin{tabular}[c]{@{}c@{}}$0.172^{***}$ \\ $(0.016)$\end{tabular} \\
		\bottomrule
	\end{tabular}
\end{table}


We believe that the observed effects may not be unique to our simulated environment. The consistent and significant shifts in discussion dynamics--such as the reduction in overall toxicity when non-minimalist strategies are applied--suggest that LLM facilitators are not inherently fixed in their behavior but rather responsive to the instructions they receive. This responsiveness is further supported by the fact that these strategies result in less heavy-handed interventions. While the dynamics of human discussions may differ, there is little reason to assume that LLM facilitators would behave in a prompt-invariant manner in real-world settings; rather, their behavior appears to be systematically influenced by the guidance provided, much like human facilitators, indicating that their actions can be intentionally shaped through appropriate instructions. Nonetheless, our results indicate that studying alternative facilitation strategies is a promising direction for future research, and experiments with human participants should be conducted to verify this finding in future work.

Finally, we note that while toxicity is a useful signal in the context of LLM-facilitated discussions, other domains may require different heuristics to detect such shifts. Thus, future work should aim to identify more general, task-agnostic methods for detecting behavioral changes among participants. This could be achieved by using reference-free LLM-as-a-judge methodologies, which could scan discussions to find changes in discussion patterns and suggest labels for subsequent human annotation~\cite{zhang2023llmaaa}, although severe limitations exist with such approaches~\cite{krumdick2025no}. 
%Unfortunately, we cannot claim that the results in Table~\ref{tab:toxicity_regression} are representative of human discussions, since we can not evaluate whether the LLM participants react to facilitation in the same ways as human discussants. In fact, there are many reasons to believe the opposite. Firstly, human participants are likely to react poorly to LLM facilitators and not treat them similarly to their human counterparts; indeed, humans react negatively to AI-generated content~\cite{alessandro2025emotional, yao2025assistants, schlesener2025understood}-although this is not certain~\cite{xu2026mitigating, schlesener2025understood}. Secondly, while our experimental assumptions have been focused on making the simulations as realistic as possible, we are still a far cry away from fully autonomous, complex, social agents. Thirdly, even if LLMs were capable of replicating the exact reactions as the `average human', there would be no robust way of evaluating this fact~\cite{rossi_2024}. However, our results clearly show that \emph{providing the LLM facilitators with different strategies changes something fundamental} in the conversational dynamics.


